\chapter{ワークスペースとパッケージ}
\label{chap:workspace}

% この章で学ぶこと
\begin{goalbox}
  \begin{itemize}
    \item ROS 2 のワークスペースの構成（\cmd{src}, \cmd{build}, \cmd{install}, \cmd{log}）
    \item colcon を使ったパッケージのビルド
    \item \cmd{ros2 pkg create} による自分のパッケージの作成
    \item \cmd{package.xml} と \cmd{setup.py}（\cmd{CMakeLists.txt}）の役割
    \item rosdep による依存関係の解決
    \item オーバーレイとアンダーレイの考え方
  \end{itemize}
\end{goalbox}

ここまでは、apt でインストールした ROS 2 のパッケージ（turtlesim など）を使ってきました。
自分でノードを書いたり、公開されているパッケージをソースコードからビルドしたりするには、\textbf{ワークスペース}という作業場所を用意します。
この章では、ワークスペースの作り方と、パッケージの作り方・ビルドの仕方を学びます。
次の\ref{chap:rclpy}章からは、ここで作ったワークスペースとパッケージを使って、Python でノードを書いていきます。

\section{ワークスペースの構成}
\label{sec:workspace-structure}

\subsection{パッケージとワークスペース}

ROS 2 では、ノードのプログラムや設定ファイルなどを、\textbf{パッケージ}という単位でまとめます（\ref{sec:what-is-ros-components}節）。
turtlesim も、\cmd{turtlesim} という名前のパッケージです。
1 つのパッケージには、関係する複数のノードや、ノードの設定ファイル、起動用のファイル（\ref{chap:launch}章）などを入れられます。

\textbf{ワークスペース}は、パッケージのソースコードを置き、ビルドするためのディレクトリです。
ワークスペースの名前は自由ですが、\cmd{\textasciitilde/ros2\_ws} とするのが一般的です。

\subsection{ワークスペースを作る}

ワークスペースを作るには、ワークスペースのディレクトリと、その中に \cmd{src} ディレクトリを作るだけです（\ref{sec:files-manipulate}節の \cmd{mkdir -p}）。

\begin{terminal}[ワークスペースを作る]
$ mkdir -p ~/ros2_ws/src
$ cd ~/ros2_ws
\end{terminal}

パッケージのソースコードは、すべて \cmd{src} ディレクトリの中に置きます。
ワークスペースでビルドを行うと、\cmd{src} と同じ階層に、自動で 3 つのディレクトリが作られます。

\begin{tablebox}[ワークスペースのディレクトリ][tab:workspace-dirs]
  \begin{tabular}{lp{0.6\linewidth}}
    \toprule
    ディレクトリ & 役割 \\
    \midrule
    \cmd{src} & パッケージのソースコードを置く。自分で編集するのはここだけ \\
    \cmd{build} & ビルドの途中で作られるファイルが置かれる \\
    \cmd{install} & ビルドしてできたものがインストールされる。ノードはここから実行される \\
    \cmd{log} & ビルドのログが置かれる \\
    \bottomrule
  \end{tabular}
\end{tablebox}

\cmd{build}・\cmd{install}・\cmd{log} は、ビルドのたびに作り直されるディレクトリです。
自分で編集する必要はなく、Git で管理するときも含めません（\ref{sec:git-basic}節の \cmd{.gitignore}）。
ビルドがおかしくなったときは、この 3 つのディレクトリを \cmd{rm -r} で削除してから、ビルドし直すと直ることがあります。

\section{colcon によるビルド}
\label{sec:workspace-colcon}

\subsection{colcon}

ROS 2 では、\textbf{colcon}（コルコン）というツールで、ワークスペースの中のパッケージをまとめてビルドします。
colcon は、\ref{sec:ros2-install-install}節でインストールした \cmd{ros-dev-tools} に含まれています。
colcon は、パッケージどうしの依存関係を調べて、正しい順番でビルドしてくれます。
C++ のパッケージは、内部で CMake（\ref{sec:package-build}節）を使ってビルドされます。

\subsection{公開されているパッケージをビルドする}

試しに、turtlesim のソースコードをダウンロードして、自分のワークスペースでビルドしてみましょう。
turtlesim のソースコードは、\cmd{ros\_tutorials} というリポジトリで公開されています。
\cmd{git clone}（\ref{sec:git-github}節）で、Jazzy 用のブランチ（\cmd{-b jazzy}）をダウンロードします。

\begin{terminal}[turtlesim のソースコードをダウンロードする]
$ cd ~/ros2_ws/src
$ git clone https://github.com/ros/ros_tutorials.git -b jazzy
\end{terminal}

ビルドする前に、パッケージが必要としている依存パッケージをインストールしておきます（詳しくは\ref{sec:workspace-rosdep}節で説明します）。

\begin{terminal}[依存パッケージをインストールする]
$ cd ~/ros2_ws
$ sudo rosdep init        # 最初の 1 回だけ
$ rosdep update
$ rosdep install -i --from-path src --rosdistro jazzy -y
#All required rosdeps installed successfully
\end{terminal}

準備ができたら、\textbf{ワークスペースのディレクトリ}（\cmd{\textasciitilde/ros2\_ws}）で \cmd{colcon build} を実行します。

\begin{terminal}[ワークスペースをビルドする]
$ cd ~/ros2_ws
$ colcon build
Starting >>> turtlesim
Finished <<< turtlesim [25.2s]

Summary: 1 package finished [25.6s]
\end{terminal}

\cmd{Summary} に、ビルドに成功したパッケージの数が表示されます。
ワークスペースの中を見ると、\cmd{build}・\cmd{install}・\cmd{log} ができています。

\begin{terminal}[ビルド後のワークスペース]
$ ls
build  install  log  src
\end{terminal}

\begin{caution}[colcon build はワークスペースのディレクトリで]
  \cmd{colcon build} は、必ずワークスペースのディレクトリ（\cmd{src} の 1 つ上）で実行します。
  \cmd{src} の中や、パッケージのディレクトリの中で実行すると、その場所に \cmd{build} などのディレクトリが作られてしまい、混乱の原因になります。
  間違えて作ってしまったら、そのディレクトリを削除しましょう。
\end{caution}

\subsection{colcon build のよく使うオプション}

\begin{tablebox}[\cmd{colcon build} のよく使うオプション][tab:workspace-colcon-options]
  \begin{tabular}{p{0.42\linewidth}p{0.5\linewidth}}
    \toprule
    オプション & 意味 \\
    \midrule
    \cmd{--packages-select パッケージ名} & 指定したパッケージだけをビルドする \\
    \cmd{--packages-up-to パッケージ名} & 指定したパッケージと、それが依存するパッケージをビルドする \\
    \cmd{--symlink-install} & ファイルをコピーせずにリンクでインストールする。Python のファイルなどを書き換えたとき、ビルドし直さずに反映される \\
    \bottomrule
  \end{tabular}
\end{tablebox}

パッケージが増えてくると、ビルドに時間がかかるようになります。
変更したパッケージだけを \cmd{--packages-select} でビルドすると、時間を短くできます。
Python のノードを書くときは、\cmd{--symlink-install} を付けておくと便利です（\ref{chap:rclpy}章）。
よく使うので、\ref{sec:env-alias}節で紹介したように、エイリアスを作っておくのもよいでしょう。

\section{パッケージの作成}
\label{sec:workspace-create}

\subsection{ros2 pkg create}

自分のパッケージは、\cmd{ros2 pkg create} コマンドで作ります。
パッケージのひな形となるディレクトリやファイルを、自動で作ってくれます。
パッケージは、ワークスペースの \cmd{src} ディレクトリの中に作ります。

\begin{terminal}[Python のパッケージを作る]
$ cd ~/ros2_ws/src
$ ros2 pkg create --build-type ament_python --license Apache-2.0 --node-name my_node my_package
going to create a new package
package name: my_package
destination directory: /home/user/ros2_ws/src
...
\end{terminal}

\begin{description}
  \item[\cmd{--build-type ament\_python}] Python でノードを書くパッケージを作ります。C++ の場合は \cmd{ament\_cmake} を指定します。
  \item[\cmd{--license Apache-2.0}] パッケージのライセンスを指定します（\ref{sec:linux-oss}節）。
  \item[\cmd{--node-name my\_node}] 指定した名前で、ノードのひな形のファイルも作ります。
  \item[\cmd{my\_package}] パッケージの名前です。パッケージ名には、小文字の英数字とアンダースコアを使います。
\end{description}

ほかのパッケージを使うノードを書く場合は、\cmd{--dependencies rclpy std\_msgs} のように、依存するパッケージも指定できます。

\subsection{パッケージの中身}

作られたパッケージの中を、\cmd{tree} コマンド（\ref{sec:package-apt}節）で見てみましょう。

\begin{terminal}[Python のパッケージの中身]
$ tree my_package
my_package
├── my_package
│   ├── __init__.py
│   └── my_node.py
├── package.xml
├── resource
│   └── my_package
├── setup.cfg
├── setup.py
└── test
    ├── test_copyright.py
    ├── test_flake8.py
    └── test_pep257.py
\end{terminal}

\begin{description}
  \item[\cmd{my\_package/}] パッケージと同じ名前のディレクトリで、ここに Python のノードのファイルを置きます。\cmd{\_\_init\_\_.py} は、このディレクトリを Python のモジュール（\ref{sec:python-class-module}節）として扱うための空のファイルです。
  \item[\cmd{package.xml}] パッケージの情報を書くファイルです（\ref{sec:workspace-files}節）。
  \item[\cmd{setup.py}, \cmd{setup.cfg}] Python のパッケージのインストール方法を書くファイルです。
  \item[\cmd{resource/}] ROS 2 がパッケージを見つけるための目印のファイルが入っています。
  \item[\cmd{test/}] コードの書き方を自動で確かめるテストが入っています。
\end{description}

\cmd{my\_node.py} には、次のようなひな形が書かれています。
\ref{sec:python-class-module}節で学んだ \cmd{if \_\_name\_\_ == '\_\_main\_\_':} が使われていることがわかります。

\begin{lstlisting}[style=python, caption={ノードのひな形（my\_package/my\_package/my\_node.py）}, label={lst:workspace-my-node}]
def main():
    print('Hi from my_package.')


if __name__ == '__main__':
    main()
\end{lstlisting}

\subsection{ビルドして実行する}

作ったパッケージをビルドして、実行してみましょう。
ビルドしたパッケージを使うには、ワークスペースの \cmd{install/local\_setup.bash} を \cmd{source} で読み込む必要があります（\ref{sec:workspace-overlay}節）。

\begin{terminal}[パッケージをビルドして実行する]
$ cd ~/ros2_ws
$ colcon build --packages-select my_package
Starting >>> my_package
Finished <<< my_package [1.21s]

Summary: 1 package finished [1.52s]
$ source install/local_setup.bash
$ ros2 run my_package my_node
Hi from my_package.
\end{terminal}

自分で作ったパッケージのノードを、\cmd{ros2 run} で実行できました。
まだ ROS 2 の通信は使っていませんが、次の\ref{chap:rclpy}章で、このパッケージにノードを書いていきます。

\section{package.xml と setup.py / CMakeLists.txt}
\label{sec:workspace-files}

\subsection{package.xml}

\cmd{package.xml} は、パッケージの名前・バージョン・説明・管理者・ライセンス、そして\textbf{依存するパッケージ}を書くファイルです。
XML（タグで囲んで書く形式）で書かれています。

\begin{lstlisting}[style=xml, caption={package.xml}, label={lst:workspace-package-xml}]
<?xml version="1.0"?>
<?xml-model href="http://download.ros.org/schema/package_format3.xsd" schematypens="http://www.w3.org/2001/XMLSchema"?>
<package format="3">
  <name>my_package</name>
  <version>0.0.0</version>
  <description>TODO: Package description</description>
  <maintainer email="user@todo.todo">user</maintainer>
  <license>Apache-2.0</license>

  <test_depend>ament_copyright</test_depend>
  <test_depend>ament_flake8</test_depend>
  <test_depend>ament_pep257</test_depend>
  <test_depend>python3-pytest</test_depend>

  <export>
    <build_type>ament_python</build_type>
  </export>
</package>
\end{lstlisting}

\cmd{description} や \cmd{maintainer} は、自分のパッケージの説明や、自分の名前とメールアドレスに書き換えておきましょう。
\cmd{license} には、\ref{sec:linux-oss}節で学んだライセンスを書きます。
パッケージを公開するときに、ほかの人が安心して使えるように、必ずライセンスを明記しておきましょう。

ほかのパッケージを使う場合は、そのパッケージ名を \cmd{<depend>} タグで書きます。
たとえば、\cmd{rclpy} と \cmd{std\_msgs} を使うノードを書くなら、次の 2 行を追加します。

\begin{lstlisting}[style=xml, caption={package.xml に依存するパッケージを追加する}, label={lst:workspace-depend}]
  <depend>rclpy</depend>
  <depend>std_msgs</depend>
\end{lstlisting}

ここに書いた依存関係をもとに、colcon はビルドの順番を決め、rosdep は必要なパッケージをインストールします（\ref{sec:workspace-rosdep}節）。

\subsection{setup.py}

Python のパッケージ（\cmd{ament\_python}）では、\cmd{setup.py} に、パッケージをどのようにインストールするかを書きます。
特に大切なのは、最後の \cmd{entry\_points} の部分です。

\begin{lstlisting}[style=python, caption={setup.py（一部）}, label={lst:workspace-setup-py}]
setup(
    name=package_name,
    ...
    license='Apache-2.0',
    entry_points={
        'console_scripts': [
            'my_node = my_package.my_node:main'
        ],
    },
)
\end{lstlisting}

\cmd{'my\_node = my\_package.my\_node:main'} は、「\cmd{ros2 run my\_package my\_node} と実行されたら、\cmd{my\_package} ディレクトリの \cmd{my\_node.py} の \cmd{main} 関数を呼び出す」という意味です。
新しいノードのファイルを追加したときは、ここに 1 行追加して、ビルドし直します（\ref{chap:rclpy}章）。

\subsection{CMakeLists.txt}

C++ のパッケージ（\cmd{ament\_cmake}）では、\cmd{setup.py} の代わりに \cmd{CMakeLists.txt}（\ref{sec:package-build}節）にビルドの方法を書きます。
C++ でノードを書く場合のほか、自分でメッセージ型を定義するパッケージ（\ref{chap:rclpy}章）でも使います。

\section{依存関係の解決}
\label{sec:workspace-rosdep}

\subsection{rosdep}

公開されているパッケージをソースコードからビルドするときは、そのパッケージが必要とするほかのパッケージ（依存パッケージ）を、先にインストールしておく必要があります。
一つ一つ調べて \cmd{apt install} するのは大変なので、ROS 2 では \textbf{rosdep} というツールを使います。

rosdep は、ワークスペースの中のパッケージの \cmd{package.xml} を読み、そこに書かれた依存パッケージを、apt のパッケージ名（\cmd{ros-jazzy-...} や \cmd{python3-...} など）に変換して、まとめてインストールしてくれます。
\ref{sec:python-pep668}節で説明したように、ROS 2 のノードで使う Python のパッケージも、rosdep を通じて apt でインストールされます。

\subsection{rosdep の使い方}

rosdep を初めて使うときは、最初に 1 回だけ \cmd{sudo rosdep init} を実行して、初期設定をします。
その後、\cmd{rosdep update} で、依存パッケージの対応表を最新にします。

\begin{terminal}[rosdep の初期設定と更新]
$ sudo rosdep init
$ rosdep update
\end{terminal}

依存パッケージをインストールするには、ワークスペースのディレクトリで次のコマンドを実行します。

\begin{terminal}[依存パッケージをインストールする]
$ cd ~/ros2_ws
$ rosdep install -i --from-path src --rosdistro jazzy -y
\end{terminal}

\begin{description}
  \item[\cmd{--from-path src}] \cmd{src} の中のすべてのパッケージの依存関係を調べる
  \item[\cmd{-i}] ワークスペースの中にあるパッケージは、インストールの対象から外す
  \item[\cmd{--rosdistro jazzy}] ROS 2 のディストリビューションを指定する
  \item[\cmd{-y}] インストールの確認に、自動で「はい」と答える
\end{description}

\begin{point}[ソースからビルドするときの手順]
  公開されているパッケージをソースコードからビルドするときは、次の流れが基本です。
  \begin{enumerate}
    \item \cmd{src} の中に \cmd{git clone} する
    \item ワークスペースのディレクトリで \cmd{rosdep install} を実行する
    \item \cmd{colcon build} でビルドする
    \item \cmd{source install/local\_setup.bash} で読み込む
  \end{enumerate}
  ただし、apt でインストールできるパッケージ（\cmd{ros-jazzy-...}）は、まず apt で入れることを考えましょう（\ref{sec:package-build}節の注意）。
\end{point}

\section{オーバーレイとアンダーレイ}
\label{sec:workspace-overlay}

\subsection{ワークスペースを重ねる}

ROS 2 では、ワークスペースを重ねて使います（図\ref{fig:workspace-overlay}）。
apt でインストールした ROS 2（\cmd{/opt/ros/jazzy}）の上に、自分のワークスペース（\cmd{\textasciitilde/ros2\_ws}）を重ねる、というイメージです。
下になる環境を\textbf{アンダーレイ}、上に重ねる環境を\textbf{オーバーレイ}といいます。

\begin{figure}[tb]
  \centering
  % 図：アンダーレイとオーバーレイ（chapters/ros2/workspace.tex）
\begin{tikzpicture}[ws/.style={draw, rounded corners=2pt, align=left, font=\small, text width=80mm, minimum height=14mm, inner sep=4pt, anchor=west}]
  \node[ws, fill=blue!8] (under) at (0,0) {\textbf{アンダーレイ}：\texttt{/opt/ros/jazzy}\\{\footnotesize apt でインストールした ROS 2（turtlesim、rclpy など）}};
  \node[ws, fill=green!10] (over) at (0,1.75) {\textbf{オーバーレイ}：\texttt{\textasciitilde/ros2\_ws/install}\\{\footnotesize 自分でビルドしたパッケージ（my\_package、改造した turtlesim など）}};
  \draw[figarrow] (8.9,-0.4) -- node[figlabel, right, align=left]{後から\\読み込む} (8.9,2.2);
  \node[figlabel, anchor=west, align=left, text=black!70] at (0,3.15) {同じ名前のパッケージがあるときは、オーバーレイ（上）のほうが使われる};
\end{tikzpicture}

  \caption{アンダーレイとオーバーレイ}
  \label{fig:workspace-overlay}
\end{figure}

ワークスペースを使うときは、まずアンダーレイ（ROS 2 本体）の設定を読み込み、その後にオーバーレイ（自分のワークスペース）の設定を読み込みます。

\begin{terminal}[アンダーレイとオーバーレイを読み込む]
$ source /opt/ros/jazzy/setup.bash          # アンダーレイ
$ source ~/ros2_ws/install/local_setup.bash # オーバーレイ
\end{terminal}

1 行目は、\ref{sec:ros2-install-setup}節で \cmd{.bashrc} に書いたので、ターミナルを開いたときに自動で実行されています。
オーバーレイを読み込むと、自分のワークスペースのパッケージも \cmd{ros2 run} などで使えるようになります。

\begin{point}[setup.bash と local\_setup.bash]
  ワークスペースの \cmd{install} には、\cmd{setup.bash} と \cmd{local\_setup.bash} の 2 つがあります。
  \cmd{local\_setup.bash} はそのワークスペースのパッケージだけを読み込み、\cmd{setup.bash} はそのワークスペースをビルドしたときのアンダーレイもまとめて読み込みます。
  アンダーレイをすでに読み込んでいる場合は、どちらを使っても同じ結果になります。
\end{point}

\subsection{オーバーレイが優先される}

アンダーレイとオーバーレイに同じ名前のパッケージがあるときは、\textbf{オーバーレイのほうが使われます}。
これを確かめるために、\ref{sec:workspace-colcon}節で自分のワークスペースにダウンロードした turtlesim を、少しだけ改造してみましょう。
turtlesim のウィンドウのタイトルを決めているのは、\cmd{turtlesim/src/turtle\_frame.cpp} の中の、次の行です。

\begin{lstlisting}[style=cpp, numbers=none, xleftmargin=0pt, framexleftmargin=0pt, caption={ros\_tutorials/turtlesim/src/turtle\_frame.cpp（一部）}, label={lst:workspace-turtle-frame}]
  setWindowTitle("TurtleSim");
\end{lstlisting}

テキストエディターや VS Code でこのファイルを開き、\cmd{"TurtleSim"} を \cmd{"MyTurtleSim"} に書き換えて保存してから、ビルドし直します。

\begin{terminal}[改造した turtlesim をビルドして起動する]
$ cd ~/ros2_ws
$ colcon build --packages-select turtlesim
$ source install/local_setup.bash
$ ros2 run turtlesim turtlesim_node
\end{terminal}

ウィンドウのタイトルが「MyTurtleSim」になっていれば、オーバーレイ（自分のワークスペース）の turtlesim が使われています。
一方、オーバーレイを読み込んでいない新しいターミナルで同じコマンドを実行すると、アンダーレイ（apt でインストールした）の turtlesim が起動し、タイトルは「TurtleSim」のままです。

このしくみを使うと、apt でインストールしたパッケージを壊すことなく、そのパッケージを自分のワークスペースで改造して試すことができます。

\begin{caution}[source を忘れずに]
  ビルドした後に \cmd{source install/local\_setup.bash} を忘れると、自分のパッケージが見つからない（\cmd{Package 'my\_package' not found}）、改造したはずなのに動作が変わらない、といったことが起こります。
  新しいターミナルを開いたときや、新しいパッケージを追加してビルドしたときは、必ず読み込み直しましょう。
  毎回入力するのが面倒なら、\cmd{.bashrc} に \cmd{source \textasciitilde/ros2\_ws/install/local\_setup.bash} を追加しておく方法もあります（ただし、ワークスペースを一度もビルドしていないとエラーになります）。
\end{caution}

\section*{この章のまとめ}

\begin{itemize}
  \item ワークスペースは、パッケージのソースコードを \cmd{src} に置いてビルドするためのディレクトリです。ビルドすると \cmd{build}・\cmd{install}・\cmd{log} が作られます。
  \item \cmd{colcon build} でワークスペースのパッケージをビルドします。必ずワークスペースのディレクトリで実行します。
  \item \cmd{ros2 pkg create} でパッケージのひな形を作れます。パッケージの情報と依存関係は \cmd{package.xml} に、Python のノードの登録は \cmd{setup.py} に書きます。
  \item rosdep を使うと、\cmd{package.xml} に書かれた依存パッケージを apt でまとめてインストールできます。
  \item apt でインストールした ROS 2（アンダーレイ）の上に、自分のワークスペース（オーバーレイ）を重ねて使います。同じ名前のパッケージは、オーバーレイが優先されます。ビルドしたら \cmd{source install/local\_setup.bash} を忘れずに。
\end{itemize}
